% Programación 3 - Ejercicio 2 (panificados)
% Compilar con:  pdflatex ejercicio2-panificados.tex
\documentclass[11pt,a4paper]{article}

\usepackage[utf8]{inputenc}
\usepackage[T1]{fontenc}
\usepackage[spanish]{babel}
\usepackage[margin=2.5cm]{geometry}
\usepackage{amsmath,amssymb,amsthm}

\theoremstyle{plain}
\newtheorem{lema}{Lema}
\newtheorem{teorema}{Teorema}

\theoremstyle{definition}
\newtheorem{definicion}{Definición}

\theoremstyle{remark}
\newtheorem*{observacion}{Observación}

\newcommand{\inv}{\operatorname{inv}}

\title{Ejercicio 2 --- Panificados}
\author{}
\date{}

\begin{document}
\maketitle
\vspace{-2.5em}

%======================================================================
\section*{Enunciado (resumen)}

Una panadería elabora $n$ panificados. Cada uno requiere una preparación de
duración $p_i$ y un horneado inmediato posterior de duración $h_i$. Las
preparaciones son independientes entre sí y pueden realizarse simultáneamente;
el horneado no, porque solo se puede hornear un panificado a la vez. Todos los
tiempos son distintos. Se busca el orden de horneado que minimiza el tiempo
total de elaboración.

%======================================================================
\section*{(a) Algoritmo}

\subsection*{Modelado}

Las preparaciones no compiten por ningún recurso: son ilimitadamente paralelas,
de modo que todas pueden comenzar en el instante $0$. Y aunque el horneado deba
ser inmediatamente posterior a la preparación, esto no agrega ninguna
restricción efectiva: siempre se puede retrasar el inicio de la preparación
para que termine justo cuando el horno queda libre.

En consecuencia, si $b_i$ denota el instante en que el panificado $i$ entra al
horno, su preparación comienza en $b_i - p_i$ y la única restricción es
\[
  b_i \;\geq\; p_i .
\]

El problema se reduce entonces a secuenciar $n$ tareas en una sola máquina (el
horno), donde la tarea $i$ tiene \emph{tiempo de liberación} $r_i = p_i$ y
\emph{duración} $h_i$, minimizando el instante en que termina la última. En
notación de scheduling, $1 \mid r_j \mid C_{\max}$.

\subsection*{Criterio}

\begin{quote}
\textbf{Hornear en orden creciente de tiempo de preparación $p_i$.}
\end{quote}

La intuición es directa tras el modelado: como todas las preparaciones arrancan
juntas, el panificado de preparación más corta es el primero en estar
disponible, y el de preparación más larga conviene hornearlo al final, porque
mientras se hornean los otros él se sigue preparando sin costo.

Obsérvese que los tiempos de horneado $h_i$ \emph{no} intervienen en el
criterio. Ordenar por $h$ decreciente, por ejemplo, es incorrecto:
\[
  n = 2, \quad p = (10,\,1), \quad h = (10,\,1)
\]
\begin{center}
\begin{tabular}{lll}
  Por $h$ decreciente: & hornea $1$ en $[10,20]$, luego $2$ en $[20,21]$
    & total $21$ \\
  Por $p$ creciente:   & hornea $2$ en $[1,2]$, \; luego $1$ en $[10,20]$
    & total $20$
\end{tabular}
\end{center}

\subsection*{Algoritmo}

\begin{verbatim}
    PANIFICADOS(p[1..n], h[1..n]):
        ordenar los índices por p ascendente
        t <- 0                        # instante en que el horno queda libre
        para cada i en ese orden:
            t <- max(t, p_i) + h_i    # espera si aún no está preparado
        devolver el orden y t
\end{verbatim}

El orden de horneado es la salida y $t$ es el tiempo total de elaboración. Si se
desea el schedule explícito, el panificado que entra al horno en $b_i$ comienza
su preparación en $b_i - p_i$.

\medskip
\noindent
\textbf{Complejidad.} $O(n\log n)$, dominada por el ordenamiento; el barrido
posterior es $O(n)$.

\medskip
\noindent
\textbf{Tiempo total.} Si $(1),(2),\dots,(n)$ denota el orden por $p$ creciente,
\[
  T \;=\; \max_{1\leq k\leq n}\;\Bigl(\, p_{(k)} + \sum_{j\geq k} h_{(j)} \Bigr).
\]

%======================================================================
\section*{(b) Demostración de corrección}

\subsection*{Formalización}

Como el horno procesa un panificado a la vez, toda solución factible determina
un orden de horneado, es decir una permutación $\sigma$ de $\{1,\dots,n\}$.
Dada $\sigma$, defino el \emph{schedule más temprano} por la recurrencia
\[
  C_0(\sigma) = 0, \qquad
  C_k(\sigma) = \max\bigl( C_{k-1}(\sigma),\, p_{\sigma(k)} \bigr) + h_{\sigma(k)},
  \qquad k = 1,\dots,n,
\]
y escribo $T(\sigma) = C_n(\sigma)$. Esto es exactamente lo que computa el
algoritmo. Una inducción inmediata muestra que cualquier schedule con el mismo
orden de horneado tiene tiempos de finalización $\geq C_k(\sigma)$ para todo
$k$; por lo tanto, para cada orden el schedule más temprano es el de menor
tiempo total, y basta comparar permutaciones. Denoto
$\mathrm{OPT} = \min_\sigma T(\sigma)$ y $\sigma_{\mathrm{alg}}$ a la
permutación que devuelve el algoritmo.

\begin{definicion}
Una \textbf{inversión} de $\sigma$ es un par de posiciones $k < l$ tal que
$p_{\sigma(k)} > p_{\sigma(l)}$; es decir, un panificado de preparación más
larga horneado antes que uno de preparación más corta. Se denota $\inv(\sigma)$
a la cantidad de inversiones de $\sigma$.
\end{definicion}

Como todos los tiempos de preparación son distintos, $\sigma_{\mathrm{alg}}$ es
la \emph{única} permutación con $\inv(\sigma) = 0$.

\subsection*{Lema de monotonía}

\begin{lema}\label{lema:mono}
Fijada una secuencia de panificados, el tiempo total de hornearla es una función
no decreciente del instante en que el horno queda libre al comenzar.
\end{lema}

\begin{proof}
Para un solo panificado de datos $(p,h)$, la función es
$g(t) = \max(t,p) + h$, que es no decreciente en $t$. Hornear una secuencia es
componer estas funciones una tras otra, y la composición de funciones no
decrecientes es no decreciente.
\end{proof}

Este lema es el que permite razonar localmente: si un cambio no aumenta el
instante en que el horno queda libre tras los panificados afectados, tampoco
aumenta el tiempo total.

\subsection*{Lema de la inversión adyacente}

\begin{lema}\label{lema:adyacente}
Si $\inv(\sigma) > 0$, entonces $\sigma$ tiene una inversión en posiciones
consecutivas.
\end{lema}

\begin{proof}
Por el contrarrecíproco. Si ninguna posición consecutiva forma inversión, vale
$p_{\sigma(1)} < p_{\sigma(2)} < \dots < p_{\sigma(n)}$, y por transitividad
$p_{\sigma(k)} < p_{\sigma(l)}$ para todo $k < l$: no hay ninguna inversión.
\end{proof}

\subsection*{Lema de intercambio}

\begin{lema}\label{lema:intercambio}
Sea $\sigma$ con una inversión en las posiciones consecutivas $k$ y $k+1$, y sea
$\sigma'$ la permutación que resulta de intercambiar esos dos panificados.
Entonces
\[
  T(\sigma') \leq T(\sigma)
  \qquad\text{y}\qquad
  \inv(\sigma') = \inv(\sigma) - 1 .
\]
\end{lema}

\begin{proof}
Sean $a = \sigma(k)$ y $b = \sigma(k+1)$, de modo que $p_a > p_b$ por ser
inversión, y sea $t = C_{k-1}(\sigma)$ el instante en que el horno queda libre
antes de la posición $k$ (con $t = 0$ si $k = 1$). Ese valor es el mismo en
$\sigma$ y en $\sigma'$, porque las posiciones anteriores no se modifican.

Aplicando la recurrencia dos veces, el horno queda libre tras las posiciones $k$
y $k+1$ en los instantes
\[
  F(\sigma)  = \max\bigl( \max(t,p_a) + h_a,\; p_b \bigr) + h_b,
  \qquad
  F(\sigma') = \max\bigl( \max(t,p_b) + h_b,\; p_a \bigr) + h_a .
\]

En la primera expresión se tiene $\max(t,p_a) + h_a \geq p_a > p_b$, de modo que
el máximo externo se alcanza en el primer argumento:
\[
  F(\sigma) = \max(t,p_a) + h_a + h_b .
\]

Para acotar $F(\sigma')$ basta ver que
$\max\bigl( \max(t,p_b) + h_b,\, p_a \bigr) \leq \max(t,p_a) + h_b$, lo cual se
verifica argumento por argumento:
\begin{itemize}
  \item $\max(t,p_b) + h_b \leq \max(t,p_a) + h_b$, porque $p_b < p_a$;
  \item $p_a \leq \max(t,p_a) \leq \max(t,p_a) + h_b$.
\end{itemize}
Por lo tanto
\[
  F(\sigma') \;\leq\; \max(t,p_a) + h_b + h_a \;=\; F(\sigma).
\]

Los panificados de las posiciones $k+2,\dots,n$ son los mismos y en el mismo
orden en ambas permutaciones; como el horno queda libre en $\sigma'$ no más
tarde que en $\sigma$, el Lema~\ref{lema:mono} da $T(\sigma') \leq T(\sigma)$.

Para el conteo de inversiones: intercambiar dos elementos \emph{consecutivos} no
altera el orden relativo de ningún otro par, ya que cada uno de los demás pares
conserva sus dos elementos en el mismo orden. El único par afectado es $(a,b)$,
que era inversión en $\sigma$ y deja de serlo en $\sigma'$. Luego
$\inv(\sigma') = \inv(\sigma) - 1$.
\end{proof}

\begin{observacion}
El Lema~\ref{lema:intercambio} es exactamente la desigualdad que justifica el
criterio de ordenamiento: conviene hornear primero al de menor tiempo de
preparación. Nótese que los tiempos de horneado $h_a$ y $h_b$ no intervienen en
la comparación.
\end{observacion}

\subsection*{Teorema}

\begin{teorema}
$T(\sigma_{\mathrm{alg}}) = \mathrm{OPT}$: el algoritmo minimiza el tiempo total
de elaboración.
\end{teorema}

\begin{proof}
Sea $\sigma^{*}$ una solución óptima, es decir
$T(\sigma^{*}) = \mathrm{OPT}$. Construyo una sucesión de permutaciones
\[
  \sigma^{*} = \sigma_0,\; \sigma_1,\; \sigma_2,\; \dots
\]
de la siguiente manera: mientras $\inv(\sigma_r) > 0$, el
Lema~\ref{lema:adyacente} garantiza que $\sigma_r$ tiene una inversión en
posiciones consecutivas, y se define $\sigma_{r+1}$ como el resultado de
intercambiarla.

Por el Lema~\ref{lema:intercambio}, en cada paso se cumple
\[
  T(\sigma_{r+1}) \leq T(\sigma_r)
  \qquad\text{e}\qquad
  \inv(\sigma_{r+1}) = \inv(\sigma_r) - 1 .
\]

Como $\inv$ es un entero no negativo que decrece exactamente en $1$ por paso, el
proceso termina tras exactamente $\inv(\sigma^{*}) \leq \binom{n}{2}$
intercambios, en una permutación $\sigma_N$ con $\inv(\sigma_N) = 0$. Siendo
todos los tiempos de preparación distintos, la única permutación sin inversiones
es la que ordena por $p$ creciente, o sea
$\sigma_N = \sigma_{\mathrm{alg}}$.

Encadenando las desigualdades de todos los pasos,
\[
  T(\sigma_{\mathrm{alg}}) = T(\sigma_N) \leq T(\sigma_{N-1}) \leq \dots
  \leq T(\sigma_0) = T(\sigma^{*}) = \mathrm{OPT},
\]
y como $\sigma_{\mathrm{alg}}$ es una solución factible vale también
$T(\sigma_{\mathrm{alg}}) \geq \mathrm{OPT}$. Por lo tanto
$T(\sigma_{\mathrm{alg}}) = \mathrm{OPT}$.
\end{proof}

%======================================================================
\section*{Comentarios sobre la redacción}

Tres puntos que conviene no omitir, porque son donde suelen perderse puntos:

\begin{itemize}
  \item \textbf{El Lema~\ref{lema:mono}.} Sin él, el Lema~\ref{lema:intercambio}
  solo prueba que el intercambio no empeora el instante en que terminan esos dos
  panificados, no el tiempo total. Es el puente entre lo local y lo global.

  \item \textbf{La terminación.} El argumento no es ``intercambio hasta que
  quede ordenado'', sino que $\inv$ decrece estrictamente y está acotada
  inferiormente por $0$, luego el proceso es finito. Eso es lo que garantiza que
  la cadena de desigualdades tenga largo finito.

  \item \textbf{La unicidad de la permutación sin inversiones.} Se apoya en que
  los $p_i$ son todos distintos, dato que el enunciado provee explícitamente. Si
  hubiera empates, $\sigma_{\mathrm{alg}}$ sería \emph{una} de las permutaciones
  sin inversiones y habría que agregar que todas ellas tienen el mismo costo.
\end{itemize}

\end{document}
